\scanpage{40}
\begin{center}-- 26 -- \hfill n° 371\end{center}
qu'il existe un homomorphisme unique de monoïdes $\varphi:L(S)\to B$ (chap. I, \uline{Rectifications} au fasc. IV) tel que $\varphi(e)=1$, $\varphi(\alpha)=\alpha$ pour tout $\alpha\in S$; par suite, si $A$ est l'algèbre du monoïde $L(S)$ par rapport à $Z$, il y a un homomorphisme unique d'anneaux $\bar\varphi:A\to B$ tel que $\bar\varphi(e)=1$ et $\bar\varphi(\alpha)=\alpha$ pour tout $\alpha\in S$. \uline{Supposons} que la relation $\bar\varphi(\lambda)=0$ implique $\lambda.x=0$ pour tout $x\in M$; autrement dit le noyau $n$ de $\bar\varphi$ est contenu dans l'annulateur de $A$-module $M$; on satisfait donc encore aux conditions 1) et 2) en considérant sur $M$ la structure de $(A/n)$-module associée à sa structure de $A$-module (chap. II, 3\textsuperscript{e} éd., § 12), et on notera que dans ce cas $A/n$ s'identifie au \uline{sous-anneau} de $B$ \uline{engendré par} $S$.

12. \uline{Exemples d'algèbres} : VI. \uline{Algèbre large d'un monoïde}.

L'algèbre d'un monoïde $S$ relative à un anneau commutatif $A$ est (en tant que $A$-module) le sous-module du produit $A^S$ formé des familles $(\alpha_s)_{s\in S}$ de support fini; la multiplication dans cette algèbre est définie par la relation $(\alpha_s)(\beta_s)=(\gamma_s)$ où, pour tout $s\in S$
\[
(25)\qquad \gamma_s=\sum_{tu=s}\alpha_t\beta_u.
\]
La somme du second membre de (25) a un sens parce que $(\alpha_s)$ et $(\beta_s)$ sont des familles de support fini, et il en est donc de même de la famille double $(\alpha_t\beta_u)_{(t,u)\in S\times S}$. Mais le second membre de (25) a encore un sens pour des éléments \uline{quelconques} $(\alpha_s)$, $(\beta_s)$ de $A^S$ lorsque le monoïde $S$ satisfait à la condition suivante :

(D) \uline{Pour tout} $s\in S$, \uline{il n'existe qu'un nombre fini de couples} $(t,u)$ \uline{d'éléments de} $S$ \uline{tels que} $tu=s$.

Supposons donc que $S$ vérifie la condition (D); sur le $A$-module produit $A^S$, on peut alors définir une loi de composition interne $((\alpha_s),(\beta_s))\to(\gamma_s)$ où, pour tout $s\in S$, $\gamma_s$ est donné par la formule (25). Il est immédiat que la multiplica-
\mspage{39}
